\begin{tabular}{llcccp{7.5cm}}
\hline
Column & Units & v10 & GAMA & SDSS & Description \\
\hline
GroupID & -- & $\checkmark$ & $\checkmark$ & $\checkmark$ & Unique group ID, as listed in the associated galaxy catalogue. GAMA: Nessie ID plus a region offset (1e5 G09, 2e5 G12, 3e5 G15, 5e5 G23) \\
Nfof & -- & $\checkmark$ & $\checkmark$ & $\checkmark$ & Group multiplicity \\
IterCenID & -- & IterCenCATAID & $\checkmark$ & $\checkmark$ & ID of the iterative central galaxy (preferred position of the group centre): GAMA III UberID for GAMA, Tempel+14 galid for SDSS. Not a CATAID \\
IterCenRA & deg & $\checkmark$ & $\checkmark$ & $\checkmark$ & RA of the iterative central galaxy (J2000) \\
IterCenDec & deg & $\checkmark$ & $\checkmark$ & $\checkmark$ & Dec of the iterative central galaxy (J2000) \\
IterCenZ & -- & $\checkmark$ & $\checkmark$ & $\checkmark$ & Redshift of the iterative central galaxy \\
Zfof & -- & $\checkmark$ & $\checkmark$ & $\checkmark$ & Median redshift of the group \\
Rad50 & $h^{-1}$\,Mpc & $\checkmark$ & $\checkmark$ & $\checkmark$ & Group radius defined by the 50th percentile group member, based on the projected comoving distance from IterCen \\
Rad1Sig & $h^{-1}$\,Mpc & $\checkmark$ & $\checkmark$ & $\checkmark$ & Group radius defined by the 68th percentile group member, based on the projected comoving distance from IterCen \\
Rad100 & $h^{-1}$\,Mpc & $\checkmark$ & $\checkmark$ & $\checkmark$ & Group radius defined by the most distant group member, based on the projected comoving distance from IterCen \\
VelDisp & km\,s$^{-1}$ & $\checkmark$ & $\checkmark$ & $\checkmark$ & Gapper velocity dispersion with VelErr subtracted in quadrature; set to zero if VelErr exceeds the raw dispersion \\
VelErr & km\,s$^{-1}$ & $\checkmark$ & $\checkmark$ & $\checkmark$ & Total group velocity error subtracted from VelDisp (input galaxy velocity errors of 50 km/s) \\
TotRmag & mag\,$-\,5\log h$ & $\checkmark$ & $\checkmark$ & $\checkmark$ & k+e corrected r-band absolute magnitude of the group, from the summed member luminosities (GAMA: DMU k+e polynomial; SDSS: Tempel+14 KCORRECT) \\
TotFluxProxy & -- & $\checkmark$ & $\checkmark$ & $\checkmark$ & Proxy for group total r-band luminosity down to M\_r - 5 log h = -14 in solar luminosities, from the observed flux and the fraction of the GAMA LF observable at the group redshift (B = 1 in Eq. 22 of Robotham et al. 2011). SDSS uses the r < 17.77 limit \\
MassProxy & $h^{-1}\,\mathrm{M}_\odot$ & $\checkmark$ & $\checkmark$ & $\checkmark$ & Dynamical mass proxy Rad50 times VelDisp squared over G, i.e. A = 1 in Eq. 18 of Robotham et al. (2011) \\
BCGID & -- & BCGCATAID & $\checkmark$ & $\checkmark$ & ID of the brightest group galaxy: GAMA III UberID for GAMA, Tempel+14 galid for SDSS. Not a CATAID \\
BCGRA & deg & $\checkmark$ & $\checkmark$ & $\checkmark$ & RA of BCG (J2000) \\
BCGDec & deg & $\checkmark$ & $\checkmark$ & $\checkmark$ & Dec of BCG (J2000) \\
CenRA & deg & $\checkmark$ & $\checkmark$ & $\checkmark$ & RA of the r-band luminosity weighted centre of the group (J2000) \\
CenDec & deg & $\checkmark$ & $\checkmark$ & $\checkmark$ & Dec of the r-band luminosity weighted centre of the group (J2000) \\
MassA & $h^{-1}\,\mathrm{M}_\odot$ & $\checkmark$ & $\checkmark$ & -- & MassProxy times the global A factor, A = 10.0 (Robotham et al. 2011 section 4.3). Calibrated on GAMA mocks: GAMA only \\
MassAfunc & $h^{-1}\,\mathrm{M}_\odot$ & $\checkmark$ & $\checkmark$ & -- & MassProxy times the functional A factor, a function of Nfof and Zfof (Robotham et al. 2011 section 4.3). Calibrated on GAMA mocks: GAMA only \\
LumB & $h^{-2}\,\mathrm{L}_\odot$ & $\checkmark$ & $\checkmark$ & -- & TotFluxProxy times the global B factor, B = 1.04 (Robotham et al. 2011 section 4.4). Calibrated on GAMA mocks: GAMA only \\
LumBfunc & $h^{-2}\,\mathrm{L}_\odot$ & $\checkmark$ & $\checkmark$ & -- & TotFluxProxy times the functional B factor, a function of Nfof and Zfof (Robotham et al. 2011 section 4.4). Calibrated on GAMA mocks: GAMA only \\
BCGZ & -- & -- & $\checkmark$ & $\checkmark$ & Redshift of BCG \\
MassHern & $h^{-1}\,\mathrm{M}_\odot$ & -- & $\checkmark$ & $\checkmark$ & Dynamical mass from Eq. 8 of Tempel et al. (2014) for a Hernquist profile: Rg = 4.582 sigma\_sky, sigma\_3D = sqrt(3) VelDisp, with sigma\_sky the projected rms member distance from IterCen. Zero where VelDisp is zero \\
MassNFW & $h^{-1}\,\mathrm{M}_\odot$ & -- & $\checkmark$ & $\checkmark$ & Dynamical mass from Eq. 8 of Tempel et al. (2014) for an NFW profile: Rg = kappa(M) sigma\_sky with kappa calibrated on the Tempel et al. (2014) NFW masses, sigma\_sky the projected rms member distance from IterCen, sigma\_3D = sqrt(3) times the rms dispersion of their Eq. 3. Zero where all members share one redshift \\
Region & -- & -- & $\checkmark$ & $\checkmark$ & Survey region: g09, g12, g15 or g23 for GAMA; SDSS for SDSS \\
\hline
\end{tabular}
